\begin{pastexam}[label={px:ntut-109-ci-03}]{北科大資工所 109 學年度 第 2 題}
（是非，只寫答案）

\textbf{True or False}. Just state the answer - you do \textbf{NOT} need to justify them. (9 pts)
\begin{enumerate}[label=(\arabic*)]
\item Let $T$ be a \textbf{minimum spanning tree} (MST) of graph $G$. Given a connected subgraph $H$
of $G$, $T\cap H$ is contained in some MST of $H$. (3 pts)
\item Are there graphs for which Prim's algorithm is faster than Kruskal's algorithm? (3 pts)
\item Dijkstra's algorithm for shortest path problem allows the edges having negative weight in the
input graph? (3 pts)
\end{enumerate}
\end{pastexam}

\begin{pastsoln}
答案：(1) True；(2) True（有）；(3) False（不允許）。

題目只要答案，下面寫出理由。

\textbf{(1) True。} $T \cap H$ 指 $T$ 裡同時也在 $H$ 裡的邊，記作 $F$。 $F$ 是 $T$ 的一部分，
沒有環。 在 $H$ 的所有最小生成樹裡，取一棵和 $F$ 共用的邊最多的，叫 $T_H$。 只要證明
$F \subseteq T_H$。

假設不是，有一條邊 $e = (u, v) \in F$ 不在 $T_H$ 裡。
\begin{itemize}
\item 把 $e$ 加進 $T_H$，在 $H$ 裡形成一個環 $C$。
\item 把 $e$ 從 $T$ 拿掉，$T$ 分成兩塊，$u$ 在一塊、$v$ 在另一塊，得到 $G$ 的一個切割。 $T$ 的邊裡
只有 $e$ 跨過這個切割。
\item 環 $C$ 從 $u$ 繞回 $v$，除了 $e$ 以外，至少還有一條邊 $f$ 跨過這個切割。 $f \ne e$，所以
$f \notin T$，當然也 $f \notin F$；$f$ 在環 $C$ 上而不是 $e$，所以 $f \in T_H$。
\item $w(e) \le w(f)$：否則 $T - e + f$ 是 $G$ 的生成樹，總權重比 $T$ 小，和 $T$ 是最小生成樹矛盾。
\item $T_H - f + e$ 仍是 $H$ 的生成樹（拿掉環上的一條邊、加上環上的另一條），總權重
$w(T_H) - w(f) + w(e) \le w(T_H)$，所以也是 $H$ 的最小生成樹。 它和 $F$ 共用的邊比 $T_H$ 多一條
（多了 $e$，拿掉的 $f$ 本來就不在 $F$ 裡），和 $T_H$ 的選法矛盾。
\end{itemize}
所以 $F = T \cap H$ 包含在 $H$ 的最小生成樹 $T_H$ 裡。 這和資結書第~24.1~節切割性質、環性質
用的是同一種交換論證。 注意結論是「某一棵」：$H$ 的最小生成樹不一定只有一棵，$T \cap H$ 也不一定
自己就是 $H$ 的生成樹（它可能不連通）。

\textbf{(2) True。} 有。 資結書第~24.3~節的比較：Kruskal 要先排序 $m$ 條邊，$O(m \lg m) = O(m \lg n)$；
Prim 用陣列存每個頂點的 key，不用堆積，是 $\Theta(n^2)$。 在稠密的圖（例如完全圖，$m = \Theta(n^2)$）
上，Kruskal 是 $\Theta(n^2 \lg n)$，Prim 陣列版是 $\Theta(n^2)$，Prim 比較快。 Prim 用費氏堆積時是
$O(m + n \lg n)$，只要 $m$ 比 $n$ 多一個超過常數的因子，例如 $m = n \lg n$，也比 Kruskal 的排序快。
反過來在稀疏的圖上兩者差不多，所以答案是「有這樣的圖」。

\textbf{(3) False。} Dijkstra 的正確性靠「從佇列取出的頂點，距離已經確定」，這要求邊的權重都不是負的
（資結書第~25.2.5~節）。 反例（就是資結書那一節的圖）：起點 $1$，邊 $1 \to 2$ 權重 $3$、$1 \to 3$ 權重 $5$、
$3 \to 2$ 權重 $-4$、$2 \to 4$ 權重 $2$，沒有負環。 Dijkstra 取出 $1$ 後 $d[2] = 3$、$d[3] = 5$；
接著取出 $2$，認定它的距離是 $3$，並由 $2 \to 4$ 得 $d[4] = 5$；剩下的 $3$、$4$ 距離都是 $5$，誰先取出都一樣。
取出 $3$、鬆弛 $3 \to 2$ 時才發現 $5 - 4 = 1 < 3$，但 $2$ 已經取出過，不會再拿 $2 \to 4$ 去更新，
$d[4]$ 停在 $5$。 真正的最短距離
是 $\delta(1, 4) = 5 - 4 + 2 = 3$（$1 \to 3 \to 2 \to 4$），答案錯了。 有負邊時要用 Bellman--Ford
（資結書第~25.3~節）。
\end{pastsoln}
